\documentclass[10pt,a4paper]{amsart}
\usepackage[margin=19mm]{geometry}
\usepackage{amsmath,amssymb,mathtools}
\usepackage[scr=rsfs,cal=euler]{mathalfa}
\usepackage{xfrac}
\usepackage[unicode,hidelinks,bookmarks=false]{hyperref}
\newcommand{\ie}{i.e.\ }
\newcommand{\cf}{cf.\ }
\newcommand{\resp}{respectively\ }
\newcommand{\Subsection}{\subsection}
\providecommand{\nobreakdash}{\nobreak\hbox{-}}
\setlength{\emergencystretch}{2em}
\multlinegap0pt
\usepackage{lmodern}
\usepackage{ifpdf}
\hypersetup{
	pdftitle={},
	pdfauthor={}
} 
\usepackage{enumitem,mathrsfs}
\usepackage{float}
\usepackage{booktabs}
\newcommand{\mincut}{{\operatorname{MinCut}}}
\usepackage{cleveref}
\usepackage{soul}
\usepackage{tikz}
\usepackage{tikz-3dplot}
\usepackage{pgfplots,pgfplotstable}
\usetikzlibrary{3d} 
\usetikzlibrary{arrows.meta}
\usepackage{algorithm}
\usepackage{algpseudocode}
\newcommand{\splt}{{\operatorname{split}}}
\newcommand{\vol}{{\operatorname{vol}}}
\newcommand{\ncut}{{\operatorname{ncut}}}
\newcommand{\supp}{{\mathsf{supp }\hspace{.075cm}}}
\newcommand{\suppsf}{{\mathsf{atom }\hspace{.075cm}}}
\newcommand{\prob}{{\mathscr{P}}}
\newcommand{\wass}{{\mathcal{W}}}
\newcommand{\setoftrees}{{\mathcal{T}}}
\newcommand{\stategraph}{{\mathbb{T}}}
\newcommand{\embedding}{{\mathsf{k}}}
\newcommand{\flowpoly}{{\mathcal{K}}}
\newcommand{\pr}[1]{\mathbb{P}\left({#1}\right)}
\newcommand{\ep}[1]{\mathbb{E}\left({#1}\right)}
\newcommand{\va}[1]{\mathrm{Var}\left({#1}\right)}
\newcommand{\ol}[1]{\overline{#1}}
\newcommand{\ord}{\mathsf{ord }\hspace{.075cm}}
\newtheorem{theorem}{Theorem}[section]
\newtheorem{corollary}[theorem]{Corollary}
\newtheorem{remark}[theorem]{Remark}
\newtheorem{definition}[theorem]{Definition}
\newtheorem{lemma}[theorem]{Lemma}
\newtheorem{proposition}[theorem]{Proposition}
\newtheorem{conjecture}[theorem]{Conjecture}
\newtheorem{assumption}[theorem]{Assumption}
\newtheorem{question}[theorem]{Question}
\newtheorem{problem}[theorem]{Problem}
\newtheorem{example}[theorem]{Example}
\usepackage{soul}
\newcommand{\hlc}[2][yellow]{{%
		\colorlet{foo}{#1}%
		\sethlcolor{foo}\hl{#2}}%
}
\newcommand{\sawyer}[1] {{ \textcolor{blue}{\hlc[blue!15]{\sf [SJR:] #1}}}}
\newcommand{\fan}[1] {{ \textcolor{red}{\hlc[red!15]{\sf [Fan:] #1}}}}
\newcommand{\attn}[1] {{ \textcolor{red}{\hlc[red!15]{\sf #1}}}}
\newcommand{\needcite}{{ \textcolor{red}{\hlc[yellow!50]{\sf [Citation needed]}}}}
\newcommand{\needdefn}{{ \textcolor{black}{\hlc[green!50]{\sf [Defn*]}}}}
\begin{document}
\title[Optimal Transport on Graphs and Stochastically Evolving Tree]{Optimal Transport on Graphs and Stochastically Evolving Trees}
\author{Fan Chung; Sawyer Jack Robertson}
\date{}
\maketitle
\noindent\emph{Source and licence.} Optimal Transport on Graphs and Stochastically Evolving Trees. Version arXiv:2608.14839v1 (CC BY 4.0). This material is distributed under the Creative Commons Attribution 4.0 International licence, http://creativecommons.org/licenses/by/4.0/, as recorded in the arXiv OAI licence metadata for this exact version and shown on the abstract page https://arxiv.org/abs/2608.14839v1. Full rights notice: licenses/arxiv-2608.14839-CC-BY-4.0.txt.\par\smallskip
\noindent\textit{Editorial scope.} Counted unit: the original \texttt{theorem} environment \texttt{thm:steepest-descent-convergence} at source lines 866--873 together with its complete proof at lines 875--885; the delivered document keeps source lines 492--919 so that every referenced label and every citation resolves inside it. Native editable source is the paper's own concatenated TeX; the AMS wrapper is editorial formatting only. No publisher class, upstream script or unrelated artwork is redistributed.
    \begin{definition}\label{def:basin}
        Let $G=(V,E)$ be a connected graph and let $g\in\mathbb{R}^{\setoftrees}$ be fixed. Let $c\in\mathbb{R}$. A nonempty subset of spanning trees $X\subseteq \setoftrees$ is called a \emph{basin of $g$ on $\stategraph$ at level $c$} if the following hold:
            \begin{enumerate}[label=(\roman*)]
                \item The induced subgraph $\stategraph[X]$ is connected,
                \item $g(T) = c$ for each $T\in X$,
                \item $g(S) > c$ for each $S\in\partial X$,
            \end{enumerate}
        where $\partial{X}$ is the vertex boundary of $X$ in $\stategraph$, defined by
            \begin{align*}
                \partial X &= \{S\in\setoftrees\setminus X \;:\; \exists\, T\in X \text{ such that } S\sim T\}.
            \end{align*}
    \end{definition}

    \begin{theorem}\label{thm:intro-basins}
        Let $G=(V,E)$ be connected. For fixed probability density vectors $\mu,\nu:V\to\mathbb{R}$, if $X\subseteq\setoftrees$ is a basin of $T\mapsto\wass_T(\mu,\nu)$ in $\stategraph$ at level $c$, then $c=\wass(\mu,\nu)$.
    \end{theorem}

    \Cref{thm:intro-basins} is proved in~\cref{sec:no-local-minima}, leading to a steepest descent algorithm for solving the optimal transport problem on $G$. Starting from an arbitrary initial tree $T_0\in\setoftrees$, we repeatedly move to a neighboring tree with strictly smaller transport cost until we arrive at a tree with cost $\wass_G(\mu,\nu)$. We show that if the input measures are assumed to have discrete values, then this procedure converges to a global minimizer in a number of steps which is polynomial in the size of the vertex set of $G$.

    \begin{theorem}\label{thm:steepest-descent-empirical}
        Let $G=(V,E)$ be a connected graph and let $\mu,\nu:V\to\mathbb{R}$ be fixed probability density vectors. Assume $\mu,\nu$ are \emph{discrete}, i.e., for some $\delta>0$, $\mu-\nu\in\delta\mathbb{Z}^{V}$. Then there exists a path of steepest descent for the transport function in $\stategraph$ arriving at the optimal value $\wass(\mu,\nu)$ in at most $\frac{n-1}{\delta}$ steps.
    \end{theorem}

    \Cref{thm:steepest-descent-empirical} is proved in~\cref{sec:steepest-descent}.

    As an application of the preceding theorems, in~\cref{sec:ollivier-ricci-curvature}, we consider the problem of computing the Ollivier--Ricci curvature of a graph. By combining our steepest descent framework with appropriate choices for the initial spanning trees, we obtain polynomial-time bounds for computing the curvature for each edge in the graph.

    The paper is organized as follows. In \cref{sec:flow-polytope} we define the tree polytope and the canonical embedding. In \cref{sec:connectivity-minimizers} we prove connectivity of global minimizers in the state graph. In \cref{sec:no-local-minima} we introduce basins and rule out suboptimal local minima. In \cref{sec:steepest-descent} we analyze steepest descent and prove a polynomial step bound. In~\cref{sec:ollivier-ricci-curvature} we apply our results to the problem of computing Ollivier--Ricci curvature on graphs. 

    \section{The tree polytope}\label{sec:flow-polytope}

    For $G=(V,E)$ fixed, we consider the associated bi-oriented graph with directed edges given by
        \begin{align}\label{eq:bi-oriented-edges}
            E' &= \{(x,y), (y,x) \,:\, \{x,y\}\in E\}.
        \end{align}
    We make use of the bi-oriented incidence matrix $B\in\mathbb{R}^{V\times E'}$, with entries given by
        \begin{align*}
            B_{v, e} &:= \begin{cases}
                +1 & \text{if } e = (v,w) \text{ for some }w\in V,\\
                -1 & \text{if } e = (w,v) \text{ for some }w\in V,\\
                0 & \text{otherwise}.
            \end{cases}
        \end{align*}
    For probability density functions $\mu,\nu:V\to\mathbb{R}$, a function $J:E'\to\mathbb{R}$ is called a $(\mu,\nu)$-feasible flow, or simply a \emph{flow}, if $J(e)\ge 0$ for each $e\in E'$ and $BJ = \mu-\nu$. The following lemma states that the optimal transport cost $\wass_G(\mu,\nu)$ can be expressed as the minimum of a linear functional over the set of all $(\mu,\nu)$-feasible flows.

    \begin{lemma}\label{lem:flow-formulation}
        Let $G=(V,E)$ be a connected graph, and let $\mu,\nu\in\prob(V)$ be fixed. Then the following holds:
            \begin{align}\label{eq:flow-formulation}
                \wass_G(\mu,\nu) &= \min\left\{\sum_{e\in E'} J(e) \,:\, J\in\mathbb{R}^{E'},\, BJ = \mu-\nu,\, J(e)\ge 0\text{ for each }e\in E'\right\}.
            \end{align}
    \end{lemma}

    \cref{lem:flow-formulation} can be proved, for example, by applying strong duality for linear programs, and we omit the proof here (see, e.g.,~\cite[Ch. 6]{peyre2019computational}, see also~\cite{essid2018quadratically}). Furthermore, in the case of a tree graph, the optimal flow achieving the value of the program~\cref{eq:flow-formulation} is uniquely determined. 

    \begin{lemma}\label{lem:transport-on-trees-formula}
        Let $T=(V,E)$ be a tree and let $\mu,\nu\in\prob(V)$ be fixed. Let $E'$ denote the bi-oriented edges of $T$ as in~\cref{eq:bi-oriented-edges}. Then the minimizer of the program~\cref{eq:flow-formulation} is unique and for each $(x,y)\in E'$ it is given by
            \begin{align}\label{eq:tree-flow-formula}
                J_T(x,y) &= \max\left\{\sum_{z\in S_x}(\mu(z)-\nu(z)),0\right\},
            \end{align}
        where $S_x$ denotes the connected component of $E\setminus\{\{x,y\}\}$ containing $x$.
    \end{lemma}

    For a proof of~\cref{lem:transport-on-trees-formula} we refer to, e.g.,~\cite{bapat1997moore} (see also~\cite{robertson2024all}). By using~\cref{eq:tree-flow-formula} specifically, we may associate to each spanning tree of $G$ a distinguished feasible flow as follows.

    \begin{definition}[Canonical embedded tree flow]\label{def:canonical-embedding}
        Let $G=(V,E)$ be connected, let $\mu,\nu\in\prob(V)$ be fixed probability density functions, and let $T\in\setoftrees$ be a fixed spanning tree of $G$. Let $T'$ denote the bi-oriented edges of $T$ as in~\cref{eq:bi-oriented-edges}. For $(x,y)\in T'$, let $S_x$ denote the connected component of $T$ containing $x$ after removing the edge $\{x,y\}$. Define the \emph{canonical embedding} of $T$ to be the vector $\embedding[T]\in \mathbb{R}_{\ge 0}^{E'}$ as follows:
            \begin{align}\label{eq:canonical-embedding}
                \embedding[T](x,y) &= \begin{cases}
                    \max\left\{\sum_{z\in S_x} (\mu(z)-\nu(z)), 0\right\}, & \text{if } \{x,y\}\in T,\\
                    0, & \text{if } \{x,y\}\notin T.
                \end{cases}
            \end{align}
    \end{definition}

    \begin{remark}\label{rmk:canonical-embedding}
        The embedding vector $\embedding[T]$ may equivalently be defined as follows. Fix a distinguished orientation of $T$ and then let $\widetilde{B}$ denote the restriction of the bi-oriented incidence matrix $B$ to the columns indexed by the oriented edges of $T$. Since $\widetilde{B}$ has trivial kernel, the linear system $\widetilde{B} X = \mu-\nu$ admits a unique solution $X$. 
        Then $\embedding[T]$ as in~\cref{eq:canonical-embedding} is recovered by setting $\embedding[T](x,y)=\max\{X(x,y),0\}$ and $\embedding[T](y,x)=\max\{-X(x,y),0\}$ for each oriented edge $(x,y)$ of $T$, and by assigning $0$ to all remaining coordinates in $E'$. 
    \end{remark}
    
    % \begin{definition}[Canonical embedded tree flow]\label{def:canonical-embedding}
    %     Let $G=(V,E)$ be connected and let $\mu,\nu\in\prob(V)$ be fixed probability density vectors. For each $T\in\setoftrees$, choose an arbitrary orientation $\widetilde{T}\subset E'$ of the edges of $T$, and let $B_T$ denote the restriction of the bi-oriented incidence matrix $B$ to the columns indexed by $\widetilde{T}$. Since $T$ is a tree, there is a unique vector $J_T\in\mathbb{R}^{\widetilde{T}}$ satisfying
    %         \begin{align}\label{eq:canonical-signed-tree-flow}
    %             B_T J_T = \mu-\nu .
    %         \end{align}
    %     We define the \emph{canonical embedded flow} of $T$ to be the vector $\embedding[T]\in\mathbb{R}^{E'}$ given, for each $(x,y)\in\widetilde{T}$, by
    %         \begin{align}\label{eq:canonical-embedding}
    %             \embedding[T](x,y) := \max\left\{J_T(x,y),0\right\},
    %             \qquad
    %             \embedding[T](y,x) := \max\left\{-J_T(x,y),0\right\},
    %         \end{align}
    %     and by $\embedding[T](e):=0$ for each $e\in E'$ of which no orientation belongs to $T$.
    % \end{definition}

    % \begin{remark}\label{rmk:canonical-embedding}
    %     \Cref{def:canonical-embedding} is independent of the auxiliary orientation $\widetilde{T}$. Indeed, reversing the orientation of a tree edge changes both the corresponding incidence column and the corresponding signed coordinate of $J_T$ by a factor of $-1$, and hence leaves the two nonnegative coordinates in~\cref{eq:canonical-embedding} unchanged. The existence and uniqueness of $J_T$ in~\cref{eq:canonical-signed-tree-flow} follow from the fact that the incidence matrix of a tree has rank $|V|-1$ and image equal to the hyperplane $\mathbf{1}^\perp$. Since $\mu$ and $\nu$ are probability densities, $\mathbf{1}^{\top}(\mu-\nu)=0$, so $\mu-\nu$ lies in this image. The vector $\embedding[T]$ is therefore a $(\mu,\nu)$-feasible flow supported on the bi-oriented edges of $T$, and along each unoriented edge ${x,y}\in T$ at most one of the two coordinates $\embedding[T](x,y)$ and $\embedding[T](y,x)$ is positive.
    % \end{remark}

    % It follows from~\cref{rmk:canonical-embedding} that the map $T\mapsto\embedding[T]:\setoftrees\to\mathbb{R}^{E'}$ is well-defined. This is the basis of the tree polytope, which we now define. If $A$ is a set of points in Euclidean space, $\mathrm{conv}(A)$ denotes its convex hull.

    \begin{definition}\label{def:flow-polytope}
        Let $G=(V,E)$ be a connected graph, and let $\mu,\nu\in\prob(V)$ be fixed. The \emph{$(\mu,\nu)$-tree polytope,} or simply the tree polytope, is given by
            \begin{align*}
                \flowpoly(\mu,\nu) &:= \mathrm{conv}\left(\{\embedding[T] \,:\, T\in\mathcal{T}(G)\} \right),
            \end{align*}
        where $\mathrm{conv}(\cdot)$ denotes the convex hull of a set in Euclidean space.
    \end{definition}    

    The advantage of using the tree polytope is that the problem of determining the transportation distance $\wass_G(\mu,\nu)$ can now be modeled on the vertices of $\flowpoly(\mu,\nu)$. As it turns out, the skeleton structure of $\flowpoly(\mu,\nu)$ is intimately related to that of the spanning tree state graph $\stategraph$. We state and prove the following proposition, which makes this precise.

    \begin{proposition}\label{prop:vertices-of-polytope}
        Let $G=(V,E)$ be a connected graph and let $\mu,\nu\in\prob(V)$ be fixed. Then the following statements hold.
            \begin{enumerate}[label=(\roman*)]
                \item For each spanning tree $T\in\setoftrees$, $\embedding[T]$ is a vertex of $\flowpoly(\mu,\nu)$.
                \item For each vertex $K\in\flowpoly(\mu,\nu)$, there exists at least one spanning tree $T\in\setoftrees$ such that $K = \embedding[T]$. Moreover, the support of $K$, i.e., edges $\{x,y\}\in E$ such that $K(x,y)>0$, forms a forest $F$ in $G$, and the fiber $\embedding^{-1}[K]$ consists of all spanning trees $T$ of $G$ containing $F$.
            \end{enumerate}
    \end{proposition}

    \begin{proof}
        For claim \emph{(i)}, fix a spanning tree $T\in\setoftrees$. Suppose $K=\embedding[T]$ is not a vertex in $\flowpoly(\mu,\nu)$. Then there exist $K_1,K_2\in\flowpoly(\mu,\nu)$ with $K_1\neq K_2$ and
            \begin{align*}
                K = \frac{1}{2}(K_1+K_2).
            \end{align*}
        In particular, $K_1,K_2\ge 0$, and thus $K(e)=0$ implies $K_1(e)=K_2(e)=0$. Let
            \begin{align*}
                \overrightarrow{T} := \{(x,y)\in E' : K(x,y) > 0\}.
            \end{align*}
        By~\cref{eq:canonical-embedding}, $\overrightarrow{T}$ is an oriented forest contained in $T$, and each of $K,K_1,K_2$ are feasible flows with supports contained in $\overrightarrow{T}$. 

        Consider the vector subspace $A\subseteq\mathbb{R}^{E'}$ given by
            \begin{align*}
                A &:= \mathrm{span} \left\{ \mathbf{1}_{(x,y)} \,:\, (x,y)\in\overrightarrow{T} \right\},
            \end{align*}
        where $\mathbf{1}_{(x,y)}$ is the standard basis vector in $\mathbb{R}^{E'}$ corresponding to the directed edge $(x,y)$. We claim that the restriction of $B$ to $A$, i.e., the linear map $B|_{A}:A\to\mathbb{R}^{V}$ given by $J\mapsto BJ$, is injective. 

        We argue by induction on $m:=|\overrightarrow{T}|$. For the base case $m=0$, we have $A=\{0\}$, so $B|_A$ is injective. For the induction step, assume $m\ge 1$ and that the claim holds for every oriented forest with fewer than $m$ edges. Let $J\in A$ satisfy $BJ=0$. Since the underlying graph of $\overrightarrow{T}$ is a nonempty forest, it contains a leaf $v\in V$ incident to a unique oriented edge $e_0\in\overrightarrow{T}$. Because $J$ vanishes outside $\overrightarrow{T}$, the $v$-th coordinate of $BJ$ is $\pm J(e_0)$, where the sign depends on the orientation of $e_0$. Hence $J(e_0)=0$. Let
            \begin{align*}
                \overrightarrow{T}_1 := \overrightarrow{T}\setminus\{e_0\}
            \end{align*}
        and set
            \begin{align*}
                A_1 := \mathrm{span}\left\{\mathbf{1}_{(x,y)}:\,(x,y)\in\overrightarrow{T}_1\right\}.
            \end{align*}
        Then $J\in A_1$, and $\overrightarrow{T}_1$ is again an oriented forest with $m-1$ edges. By the induction hypothesis, the restriction of $B$ to $A_1$ is injective, so $BJ=0$ implies $J=0$. This proves that $B|_A$ is injective.

        Since $K,K_1,K_2\in A$ and each is a feasible flow, we have
            \begin{align*}
                BK = BK_1 = BK_2 = \mu-\nu.
            \end{align*}
        Injectivity of $B|_A$ therefore gives $K=K_1=K_2$, contradicting $K_1\neq K_2$. This shows $\embedding[T]$ is a vertex of the polytope $\flowpoly(\mu,\nu)$, proving claim \emph{(i)}.

        For claim \emph{(ii)}, since $\flowpoly(\mu,\nu)$ is the convex hull of $\{\embedding[T]:T\in\setoftrees\}$, a vertex $K$ of $\flowpoly(\mu,\nu)$ is necessarily one of these points and hence $K=\embedding[T]$ for some spanning tree $T\in\setoftrees$. Let $F = \{\{x,y\}\in E : K(x,y) > 0\}$. By~\cref{eq:canonical-embedding}, the set $F$ is contained in $T$, hence $F$ is a forest.

        We next show that
            \begin{align*}
                \embedding^{-1}[K] = \{S\in\setoftrees : S\supseteq F\}.
            \end{align*}
        First, if $\embedding[S]=K$, then every positive coordinate of $K$ is supported on an edge of $S$, so necessarily $S\supseteq F$.

        Conversely, let $S\in\setoftrees$ satisfy $S\supseteq F$. Orient the edges of $S$ so that each edge of $F$ points in the direction where $K$ is positive, and orient the remaining edges arbitrarily. Now assign each edge of $F$ the corresponding value of $K$, and assign each edge of $S\setminus F$ the value $0$. This is just the flow $K$ written in the chosen orientation, so it is feasible and its positive part is exactly $K$. By the uniqueness statement in~\cref{rmk:canonical-embedding}, this is the signed flow associated with $S$, and therefore $\embedding[S]=K$. This proves the reverse direction and completes the proof of claim \emph{(ii)}.
    \end{proof}

\section{Properties of the tree polytope}\label{sec:connectivity-minimizers}

    In this section, we use the tree polytope to prove that the global minimizers of the function $T\mapsto \wass_T(\mu,\nu)$ form a connected set in the state graph $\stategraph$. To do so, we state and prove a lemma about the connectivity structure of the state graph.

    \begin{proposition}\label{prop:connected-fibers}
        Let $G=(V,E)$ be a connected graph and let $F\subseteq E$ be a forest in $G$. Let
            \begin{align*}
                \setoftrees[F] &= \{T\in\setoftrees : T\supseteq F\}.
            \end{align*}
        Then the induced subgraph $\stategraph[\setoftrees[F]]$ of the spanning tree state graph $\stategraph$ is connected.
    \end{proposition}

    \begin{proof}
        If $|F|=|V|-1$, then $F$ is already a spanning tree and $\setoftrees[F]=\{F\}$, so the claim is trivial. Assume $|F|=k<|V|-1$. We will show that for any $T,T'\in\setoftrees[F]$, there is a path in $\stategraph[\setoftrees[F]]$ connecting $T$ and $T'$.

        If $T=T'$, there is nothing to prove. Otherwise choose an edge $e\in T\setminus T'$. Set $F_1 := T\setminus\{e\}$. Then $F_1$ is a forest with two connected components on vertex sets $A_1$ and $B_1$. Since $T'$ is connected, it must contain at least one edge in the cut $(A_1,B_1)$. Because $e$ is the unique edge of $T$ crossing this cut, any such edge lies in $T'\setminus T$. Choose one such edge $f\in E(A_1,B_1)\cap T'$. Then
            \begin{align*}
                T_1 := T\setminus\{e\}\cup\{f\}
            \end{align*}
        is a spanning tree, and we note $T_1$ still contains $F$ because $e\notin F$. 
        Moreover, since $e\notin T'$ and $f\in T'$, we have $T_1\cap T' = (T\cap T')\cup\{f\}$, and therefore $|T_1\cap T'| = |T\cap T'|+1$. Repeating this process increases the number of edges shared with $T'$ by one at each step. Since a spanning tree has only $|V|-1$ edges, after finitely many steps we must arrive at $T'$, producing a path from $T$ to $T'$ in $\stategraph[\setoftrees[F]]$.
    \end{proof}

    Next, we show that edges in $\flowpoly(\mu,\nu)$ are lifted in a straightforward way to links in $\stategraph$, as follows.

    \begin{lemma}\label{lem:adjacency-in-K}
        Let $G=(V,E)$ be a connected graph and let $\mu,\nu\in\prob(V)$ be fixed. Let $K_1, K_2\in\flowpoly(\mu,\nu)$ be vertices such that $\{K_1, K_2\}$ is an edge ($1$-face) of $\flowpoly(\mu,\nu)$. If $T \in\setoftrees$ satisfies $\embedding[T] = K_1$, then there exists $T'\in\setoftrees$ such that $\embedding[T'] = K_2$ and such that $T$ and $T'$ are adjacent in the state graph $\stategraph$.
    \end{lemma}

    \begin{proof}
        Let $K_1, K_2\in\flowpoly(\mu,\nu)$ be vertices such that $\{K_1,K_2\}$ is an edge of $\flowpoly(\mu,\nu)$, and let $T\in\setoftrees$ satisfy $\embedding[T]=K_1$. By~\cref{prop:vertices-of-polytope}(ii), there exist forests $F_1,F_2\subseteq E$, obtained as the supports of $K_1$ and $K_2$, respectively, and with the property that
            \begin{align}\label{eq:fiber-description-adjacency}
                \embedding^{-1}[K_i] = \{S\in\setoftrees: S\supseteq F_i\},\text{ for }\quad i=1,2.
            \end{align}
        In particular, $T\supseteq F_1$. Since $\{K_1,K_2\}$ is an edge, we have $K_1\neq K_2$. Hence the fibers in~\cref{eq:fiber-description-adjacency} are disjoint and in particular there is no spanning tree containing both $F_1$ and $F_2$, so the union $F = F_1\cup F_2$ cannot be a forest. Thus, by passing to a simple subcycle by deleting repeated vertices as needed, $F$ contains at least one simple cycle. We claim that this cycle is unique. Suppose $F$ contains two distinct simple cycles. Let $\overline{K} := \tfrac{1}{2}(K_1+K_2)$ and set
            \begin{align*}
                F^+ := \{(x,y)\in E' : \overline{K}(x,y) > 0\}.
            \end{align*}
        Then $\overline{K}$ is strictly positive on $F^+$ and satisfies $B_{F^+}\overline{K}^{(F^+)} = \mu-\nu$, where $\overline{K}^{(F^+)}$ denotes the restriction of $\overline{K}$ to the coordinates indexed by $F^+$, and where $B_{F^+}$ denotes the restriction of $B$ to the columns indexed by $F^+$. The underlying undirected graph of $F^+$ contains $F$, so if $F$ has two distinct cycles then $\dim\ker(B_{F^+})\ge 2$. Choose linearly independent $h_1,h_2\in\ker(B_{F^+})$. For $|\varepsilon_1|,|\varepsilon_2|$ sufficiently small, the vector
            \begin{align*}
                \widetilde{K} = \overline{K} + \varepsilon_1 h_1 + \varepsilon_2 h_2
            \end{align*}
        remains nonnegative and satisfies $B\widetilde{K} = \mu-\nu$, so $\widetilde{K}\in\flowpoly(\mu,\nu)$. This produces a $2$-dimensional face through $\overline{K}$, contradicting that $\{K_1,K_2\}$ is an edge. Therefore $F$ contains a unique simple cycle, say $C\subseteq F$.
        
        Next we claim that $F_1$ and $F_2$ differ only on $C$; specifically,
            \begin{align*}
                F_1\cap(F\setminus C) = F_2\cap(F\setminus C) = F\setminus C.
            \end{align*}
        To see this, let $g\in F\setminus C$ be an edge off the cycle. We claim that $g\in F_1\cap F_2$. Clearly $g\in F_1$ or $g\in F_2$ holds automatically, so assume without loss of generality that $g\in F_1$. If $g\notin F_2$, then $K_2$ vanishes on both orientations of $g$ while $K_1$ is positive on exactly one orientation $\overrightarrow{g}$ of $g$, so $K_1-K_2$ is nonzero on $\overrightarrow{g}$. Since $B(K_1-K_2)=0$, $F$ contains a unique cycle, so the support of $K_1-K_2$ must contain $C$. This is impossible because $g\notin C$. Therefore $g\in F_2$, and we conclude that $F_1\cap(F\setminus C) = F_2\cap(F\setminus C) = F\setminus C$.

        We claim that $C$ contains at least one edge $e\in F_2$ such that $e\notin T$. Indeed, if no such edge existed, then $C\subseteq F_1\subseteq T$, contradicting that $T$ is a tree. Therefore $C$ must contain an edge avoiding $T$. Since $T$ contains $F_1$ this edge must lie in $F_2$. Fix such an edge $e$. Adding $e$ to $T$ creates a unique simple cycle; denote it by $C_T(e)\subseteq T\cup\{e\}$. Indeed, $C_T(e)$ is exactly the edge $e$ together with the unique simple path in $T$ joining the endpoints of $e$. Every edge of $C_T(e)$ other than $e$ lies in $T$, and $C_T(e)$ must contain at least one edge of $T$ that is not in $F_2$, otherwise $C_T(e)\subseteq (T\cup\{e\})\cap F_2$ would contradict that $F_2$ is a forest. Fix an edge $e'\in C_T(e) \cap (T\setminus F_2)$. Define
            \begin{align*}
                T' \;:=\; (T\cup\{e\})\setminus\{e'\}.
            \end{align*}
        Clearly $T'\in\setoftrees$, and by construction, $T\sim T'$ in $\stategraph$. By~\eqref{eq:fiber-description-adjacency}, it remains to show that $T'\supseteq F_2$. Let $f\in F_2$ be arbitrary. If $f=e$, then $f\in T'$ by construction. Thus we may assume $f\neq e$ and consider two cases.
        
        First, if $f\notin C$, then $f\in F_2\cap(F\setminus C)=F\setminus C$ by construction, hence $f\in F_1$ and therefore $f\in T$ (since $T\supseteq F_1$). Moreover, to produce $T'$ we removed only $e'\in C$, so $f$ remains in $T'$. Thus $f\in T'$.
        
        Second, if $f\in C$, then the only edge of $C$ removed in passing from $T\cup\{e\}$ to $T'$ is $e'$, and by construction $e'\notin F_2$. But $f\in F_2$ so $f\neq e'$ and hence $f\in T'$.

        In both cases we have shown $f\in T'$. Since $f\in F_2$ was arbitrarily chosen, this proves $T'\supseteq F_2$. By~\cref{prop:vertices-of-polytope}(ii), it follows that $\embedding[T']=K_2$. Therefore we have proved that $T\sim T'$ in $\stategraph$, $\embedding[T]=K_1$, and $\embedding[T']=K_2$ as desired.
    \end{proof}

    Finally, using the preceding setup, we will show that the minimizers of the transport function are connected.

    \begin{theorem}\label{thm:connectivity-minimizers}
        Let $G=(V,E)$ be a connected graph and let $\mu,\nu\in\prob(V)$ be fixed. Consider the minimum level set $X = \{T\in\setoftrees : \wass_{T}(\mu,\nu) = \wass(\mu,\nu)\}$ of the transport function $T\mapsto \wass_{T}(\mu,\nu)$. Then the induced subgraph $\stategraph[X]$ of the spanning tree state graph $\stategraph$ is connected.
    \end{theorem}

    \begin{proof}
        Consider the $(\mu,\nu)$-tree polytope $\flowpoly(\mu,\nu)$, and let
            \begin{align*}
                \flowpoly_\star &= \flowpoly(\mu,\nu) \cap \{K\in\mathbb{R}^{E'} : \mathbf{1}^\top K = \wass(\mu,\nu)\}.
            \end{align*}
        We claim that $\flowpoly_\star$ is a face of $\flowpoly(\mu,\nu)$. Recall a basic fact from polytopal geometry: if $F$ is a subset of a convex polytope $P\subseteq\mathbb{R}^{d}$, then $F$ is a face of $P$ if and only if it is convex and satisfies the following condition: if $x\in F$ and $x = \theta y + (1-\theta) z$ with $y, z\in P$ and $0\le \theta\le 1$, then $y, z\in F$. To this end, let $K\in\flowpoly_\star$, $K_1, K_2\in\flowpoly(\mu,\nu)$, and $0 \le \theta \le 1$ be fixed, and assume $K = \theta K_1 + (1-\theta) K_2$. Then
            \begin{align*}
                \mathbf{1}^\top K &= \theta \mathbf{1}^\top K_1 + (1-\theta) \mathbf{1}^\top K_2.
            \end{align*}
        Since $\flowpoly(\mu,\nu)$ is the convex hull of $\{\embedding[T]:T\in\setoftrees\}$, we have $\mathbf{1}^\top K_i \ge \wass(\mu,\nu)$ for $i=1,2$. Together with $\mathbf{1}^\top K = \wass(\mu,\nu)$ and the fact that $K(e)\geq 0$ for each $e\in E'$, it follows that $\mathbf{1}^\top K_i = \wass(\mu,\nu)$ for $i=1,2$, and hence that $K_i\in\flowpoly_\star$. This proves that $\flowpoly_\star$ is a face of $\flowpoly(\mu,\nu)$.

        Now take two trees $S, T\in \setoftrees$ achieving the optimal transport cost between $\mu,\nu$, i.e., $S,T\in X$. We claim that there exists a path in $\stategraph[X]$ connecting them. Clearly $\embedding[S], \embedding[T]\in \flowpoly_\star$. Since $\flowpoly_\star$ is a nonempty face of $\flowpoly(\mu,\nu)$, its $1$-skeleton is connected, so there exists a path
            \begin{align*}
                (K_0 = \embedding[S], K_1,\dots, K_{r-1}, K_{r} = \embedding[T])
            \end{align*}
        of vertices in $\flowpoly_\star$ connecting the flows $\embedding[S], \embedding[T]$. Since $\flowpoly_\star$ is a face of $\flowpoly(\mu,\nu)$, these are also vertices of $\flowpoly(\mu,\nu)$.

        Starting from $T_0 = S$, apply~\cref{lem:adjacency-in-K} to each edge $\{K_{i-1},K_i\}$ to obtain $T_i\in\setoftrees$ with $\embedding[T_i]=K_i$ and $T_{i-1}\sim T_i$ in $\stategraph$. Because $K_i\in\flowpoly_\star$, we have $\mathbf{1}^\top \embedding[T_i] = \wass(\mu,\nu)$, and hence $T_i\in X$. This yields a path
            \begin{align*}
                (T_0 = S, T_1, \dotsc, T_{r-1}, T_{r})
            \end{align*}
        in $\stategraph[X]$ with $\embedding[T_i]=K_i$ for each $0\le i\le r$.

        The endpoint $T_{r}$ need not equal $T$, but they belong to the same fiber of $\embedding$ since $\embedding[T]=\embedding[T_r]$. By~\cref{prop:connected-fibers}, this fiber is connected in $\stategraph$, so we may adjoin a path from $T_{r}$ to $T$. All vertices on this additional path share the same value of $\mathbf{1}^\top \embedding[\cdot]$, hence remain in $X$. The proof is completed.
    \end{proof}

\section{Properties of minimizers of the transport function}\label{sec:no-local-minima}

Before we proceed to show that the function $T\mapsto \wass_T(\mu,\nu)$ on the spanning tree state graph $\stategraph$ has no suboptimal local minima, we first state and prove a useful lemma.

\begin{lemma}[Edge descent at nonoptimal vertices]\label{lem:edge-descent-polytope}
    Let $K_0$ be a vertex of $\flowpoly(\mu,\nu)$ such that $\mathbf{1}^\top K_0 > \wass(\mu,\nu)$. Then there exists a vertex $K_1$ adjacent to $K_0$ in the $1$-skeleton of $\flowpoly(\mu,\nu)$ with
        \begin{align*}
            \mathbf{1}^\top K_1 < \mathbf{1}^\top K_0.
        \end{align*}
\end{lemma}

\begin{proof}
    Let $K_\star\in\flowpoly(\mu,\nu)$ be optimal for the linear objective $K\mapsto \mathbf{1}^\top K$, so that $\mathbf{1}^\top K_\star = \wass(\mu,\nu)$. Set $d:=K_\star-K_0$. Since $\flowpoly(\mu,\nu)$ is convex, we have
        \begin{align*}
            K_0+t d = (1-t)K_0+tK_\star \in \flowpoly(\mu,\nu),\quad t\in[0,1],
        \end{align*}
    so $d$ is a feasible direction at $K_0$.

    For each vertex $K_1',K_2',\dotsc, K_m'$ of $\flowpoly(\mu,\nu)$ adjacent to $K_0$ in its $1$-skeleton, let $d^{(i)}$ denote the edge direction $d^{(i)} = K_{i}'-K_0$, $i=1,2,\dotsc, m$. A standard polytope fact is that the feasible cone at a vertex is generated by its incident edge directions (see, e.g.,~\cite{ziegler1995lectures}), hence there are coefficients $\alpha_i\ge 0$ such that
        \begin{align*}
            d = \sum_{i=1}^m \alpha_i d^{(i)}.
        \end{align*}
    Taking objective inner products gives
        \begin{align*}
            0 > \mathbf{1}^\top d = \mathbf{1}^\top K_\star - \mathbf{1}^\top K_0 = \sum_{i=1}^m \alpha_i\,\mathbf{1}^\top d^{(i)}.
        \end{align*}
    Therefore $\mathbf{1}^\top d^{(i_0)}<0$ for some $i_0$. Let $K_1$ be the other endpoint of the corresponding edge through $K_0$; then $K_1 = K_0 + \tau d^{(i_0)}$ for some $\tau>0$, and
        \begin{align*}
            \mathbf{1}^\top K_1 = \mathbf{1}^\top K_0 + \tau\,\mathbf{1}^\top d^{(i_0)} < \mathbf{1}^\top K_0.
        \end{align*}
\end{proof}

The next result rules out isolated \emph{vertices} of the state graph $\stategraph$ where the transport function might have a strict suboptimal local minimum.

\begin{theorem}\label{thm:isolated-vertex-minimizers}
    Let $G=(V,E)$ be connected, $\mu,\nu\in\prob(V)$, and let $\setoftrees$ denote the set of spanning trees of $G$. Consider the function $f\in\mathbb{R}^{\setoftrees}$ defined by
        \begin{align*}
            f(T) &= \wass_T(\mu,\nu), \quad T\in\setoftrees.
        \end{align*}
    Let $T_0\in\setoftrees$ be fixed. Assume $T_0$ is a strict minimizer of $f$ on its neighborhood in $\stategraph$, i.e.,
        \begin{align}\label{eq:strict-local-min}
            f(S) > f(T_0)\text{ for each }S\sim T_0.
        \end{align}
    Then $T_0$ is a global minimizer of $f$, i.e.,
        \begin{align*}
            f(T_0) = \wass(\mu,\nu).
        \end{align*}
\end{theorem}

% \begin{proof}
%     Let $K_0 := \embedding[T_0]\in\flowpoly(\mu,\nu)$. By construction of the canonical embedding,
%         \begin{align*}
%             f(T_0) = \wass_{T_0}(\mu,\nu) = \mathbf{1}^\top K_0.
%         \end{align*}
%     Suppose, toward a contradiction, that $f(T_0)>\wass(\mu,\nu)$. Then $\mathbf{1}^\top K_0>\wass(\mu,\nu)$, so by~\cref{lem:edge-descent-polytope} there exists a vertex $K_1\in\flowpoly(\mu,\nu)$ adjacent to $K_0$ with
%         \begin{align*}
%             \mathbf{1}^\top K_1 < \mathbf{1}^\top K_0.
%         \end{align*}
%     By~\cref{lem:adjacency-in-K}, there exists $S\in\setoftrees$ such that $S\sim T_0$ in $\stategraph$ and $\embedding[S]=K_1$. Using again the canonical embedding identity $f(T)=\mathbf{1}^\top\embedding[T]$, we obtain
%         \begin{align*}
%             f(S) = \mathbf{1}^\top K_1 < \mathbf{1}^\top K_0 = f(T_0),
%         \end{align*}
%     contradicting~\cref{eq:strict-local-min}. Hence $f(T_0)=\wass(\mu,\nu)$.
% \end{proof}
 
In fact, a stronger version of~\cref{thm:isolated-vertex-minimizers} holds with vertices replaced by ``basins,'' as defined in~\cref{def:basin}. 

\begin{theorem}\label{thm:no-isolated-basins}
    Let $G=(V,E)$ be a connected graph, $\mu,\nu\in\prob(V)$ be fixed, let $\setoftrees$ denote the set of spanning trees of $G$, and let $\stategraph$ denote the spanning tree state graph. Let $f:\setoftrees\rightarrow\mathbb{R}$ denote the transport function $f(T) = \wass_T(\mu,\nu)$. Let $X\subseteq\setoftrees$ be a basin of $f$ at level $c\in\mathbb{R}$. Then $c = \wass_G(\mu,\nu)$.
\end{theorem}

\begin{proof}[Proof of~\cref{thm:no-isolated-basins}]
    Suppose to the contrary that $c>\wass(\mu,\nu)$. Fix any $T\in X$ and set $K_0:=\embedding[T]$. Then
        \begin{align*}
            \mathbf{1}^\top K_0 = f(T)=c>\wass(\mu,\nu).
        \end{align*}
    By~\cref{lem:edge-descent-polytope}, there exists a vertex $K_1\in\flowpoly(\mu,\nu)$ adjacent to $K_0$ with
        \begin{align*}
            \mathbf{1}^\top K_1 < \mathbf{1}^\top K_0 = c.
        \end{align*}
    Applying~\cref{lem:adjacency-in-K} to the edge $\{K_0,K_1\}$, we obtain $S\in\setoftrees$ such that $S\sim T$ in $\stategraph$ and $\embedding[S]=K_1$. Hence
        \begin{align*}
            f(S)=\mathbf{1}^\top K_1 < c.
        \end{align*}
    But $S\sim T$ with $T\in X$, so either $S\in X$ or $S\in\partial X$. If $S\in X$, then by the definition of basin we must have $f(S)=c$, leading to a contradiction. If $S\in\partial X$, then again by definition we must have $f(S)>c$, similarly leading to a contradiction. Therefore $c\le \wass(\mu,\nu)$.

    Since $\wass(\mu,\nu)$ is the global minimum of $f$, we always have $c\ge \wass(\mu,\nu)$ for any nonempty level set $\{T: f(T)=c\}$. Hence $c=\wass(\mu,\nu)$.
\end{proof}

\section{Convergence of steepest descent on the state graph}\label{sec:steepest-descent}

In this section we analyze a simple steepest descent algorithm for minimizing the transport function $f(T) = \wass_T(\mu,\nu)$ on the spanning tree state graph $\stategraph$. Before defining the algorithm, we state and prove a technical lemma.

\begin{lemma}\label{lem:descent-step}
    Let $G=(V,E)$ be a connected graph and let $\mu,\nu\in\prob(V)$ be fixed. Let $T\in\setoftrees$ be fixed satisfying
        \begin{align*}
            \wass_T(\mu,\nu) > \wass(\mu,\nu).
        \end{align*}
    Then there exists a tree $T'\in\setoftrees$ with $T'\sim T$ in $\stategraph$ such that
        \begin{align*}
            \wass_{T'}(\mu,\nu) < \wass_T(\mu,\nu).
        \end{align*}
\end{lemma}

\begin{proof}
    Set $K_0:=\embedding[T]$. By~\cref{prop:vertices-of-polytope}, $K_0$ is a vertex of $\flowpoly(\mu,\nu)$. By assumption it follows that
        \begin{align*}
            \mathbf{1}^{\top}K_0 > \wass(\mu,\nu).
        \end{align*}
    Applying~\cref{lem:edge-descent-polytope}, there exists a vertex $K_1$ adjacent to $K_0$ in the $1$-skeleton of $\flowpoly(\mu,\nu)$ such that
        \begin{align*}
            \mathbf{1}^{\top}K_1 < \mathbf{1}^{\top}K_0.
        \end{align*}
    Since $K_0=\embedding[T]$ and $\{K_0,K_1\}$ is an edge of $\flowpoly(\mu,\nu)$, \cref{lem:adjacency-in-K} yields a spanning tree $T'\in\setoftrees$ such that $T'\sim T$ and $\embedding[T']=K_1$. By construction, $\mathbf{1}^{\top}\embedding[S]=\wass_S(\mu,\nu)$ for every spanning tree $S$, so it follows that $\wass_{T'}(\mu,\nu)=\mathbf{1}^{\top}K_1<\mathbf{1}^{\top}K_0=\wass_T(\mu,\nu)$. This proves the lemma.
\end{proof}

Using~\cref{lem:descent-step}, the steepest descent algorithm on $\stategraph$ proceeds as follows. Starting from an initial tree $T_0\in\setoftrees$, if no neighbors of $T_0$ have strictly lower cost, then $T_0$ is optimal and we are done. Otherwise there exists a neighbor $T_1\sim T_0$ with $\wass_{T_1}(\mu,\nu) < \wass_{T_0}(\mu,\nu)$. We repeat this process by checking the neighbors of $T_1$. Since the state graph is finite, this process terminates in finitely many steps with an optimal tree. 


\begin{theorem}\label{thm:steepest-descent-convergence}
    Let $G=(V,E)$ be a connected graph and let $\mu,\nu\in\prob(V)$ be fixed with $\mu\ne \nu$. Let $\delta = \delta(\mu,\nu) > 0$ denote the smallest positive improvement gap between two trees with differing transport cost. Then for any initial tree $T\in\setoftrees$ there exists a path $T_0=T,T_1,\dotsc, T_{m}$ in $\stategraph$ satisfying:
        \begin{enumerate}
            \item All steps lead to strict improvement of the transport cost, i.e., $\wass_{T_{i+1}}(\mu,\nu) < \wass_{T_i}(\mu,\nu)$ for each $0\le i \le m-1$;
            \item The length of the path satisfies $m\leq \frac{n-1}{\delta}$;
            \item The path terminates with optimal cost, i.e., $\wass_{T_m}(\mu,\nu) = \wass(\mu,\nu)$.
        \end{enumerate}
\end{theorem}

\begin{proof}
    If $T$ is optimal or if $\delta = +\infty$, the claim holds immediately. We may assume $\wass_{T}(\mu,\nu) > \wass(\mu,\nu)$. 
    
    By~\cref{lem:descent-step}, there exists a neighbor $T_1\sim T$ in $\stategraph$ with $\wass_{T_1}(\mu,\nu) < \wass_{T}(\mu,\nu)$. If $\wass_{T_1}(\mu,\nu)=\wass(\mu,\nu)$, we are done; otherwise at stage $i$, having constructed $T_i$ with $\wass_{T_i}(\mu,\nu) > \wass(\mu,\nu)$, we apply the lemma to obtain $T_{i+1}\sim T_i$ with $\wass_{T_{i+1}}(\mu,\nu) < \wass_{T_i}(\mu,\nu)$.

    This greedy descent process terminates in finitely many steps since each step yields a strict improvement of at least $\delta > 0$. Specifically, starting from $\wass_{T}(\mu,\nu)$, each step decreases the cost by at least $\delta$, so the number of steps is bounded by
        \begin{align*}
            m \leq \frac{\wass_{T}(\mu,\nu) - \wass(\mu,\nu)}{\delta} \leq \frac{\max_{T'\in\setoftrees} \wass_{T'}(\mu,\nu) - \wass(\mu,\nu)}{\delta} \leq \frac{n-1}{\delta},
        \end{align*}
    where the last inequality follows because every spanning tree on $n$ vertices has $n-1$ edges and for each oriented edge $(x,y)$ it follows that the optimal flow $J_{T}$ as in~\cref{eq:tree-flow-formula} satisfies $J_T(x,y)\le 1$. Therefore $\wass_{T'}(\mu,\nu)\le n-1$ for each $T'\in\setoftrees$. By construction, the path $T_0, T_1, \dotsc, T_m$ satisfies all three required properties.
\end{proof}


The minimum improvement gap $\delta$ may be arbitrarily small in general, but the following result gives a nontrivial lower bound in the case where $\mu$ and $\nu$ are discrete measures, which arise in various applications.

\begin{lemma}\label{lem:rational-measures-gap}
    Let $G=(V,E)$ be a connected graph and let $\mu,\nu\in\prob(V)$ be fixed with $\mu\ne \nu$. Assume that for each $x\in V$ it holds $\mu(x),\nu(x) \in \alpha\mathbb{Z}$ for some $\alpha > 0$. Then $\delta(\mu,\nu) \geq \alpha$.
\end{lemma}

\Cref{lem:rational-measures-gap} is a simple consequence of~\cref{lem:transport-on-trees-formula}: for each spanning tree $T$ of $G$, $\wass_T(\mu,\nu)$ is a nonnegative integer multiple of $\alpha$, and thus any two distinct tree transport values differ by at least $\alpha$. 

% \begin{proof}
%     Let $b = \mu - \nu \in \mathbb{R}^V$. By assumption, $b(x) \in \alpha\mathbb{Z}$ for each $x \in V$, and $\mathbf{1}^\top b = 0$ holds since $\mu,\nu\in\prob(V)$. Consider two distinct spanning trees $T, T' \in \setoftrees$ with $\wass_T(\mu, \nu) \neq \wass_{T'}(\mu, \nu)$. 
%     
%     Choose arbitrary orientations of the edges of $T$ and $T'$. By~\cref{lem:transport-on-trees-formula}, the corresponding optimal tree flows satisfy
%         \begin{align*}
%             J_T(x,y) &= \max\left\{\sum_{z\in S_x} b(z),0\right\},
%             \qquad
%             J_{T'}(u,v) = \max\left\{\sum_{z\in S_u} b(z),0\right\},
%         \end{align*}
%     for each oriented edge $(x,y)$ of $T$ and $(u,v)$ of $T'$, where $S_x$ and $S_u$ denote the connected components on the tail sides after deleting the corresponding edge. Since each sum $\sum_{z\in S_x} b(z)$ and $\sum_{z\in S_u} b(z)$ lies in $\alpha\mathbb{Z}$, every such flow value lies in $\alpha\mathbb{Z}_{\ge 0}$, and every nonzero one is at least $\alpha$.
%
%     Thus both $\wass_T(\mu, \nu)$ and $\wass_{T'}(\mu, \nu)$ are nonnegative integer multiples of $\alpha$. Since $\wass_T(\mu, \nu) \neq \wass_{T'}(\mu, \nu)$, we have
%         \begin{align*}
%             |\wass_T(\mu, \nu) - \wass_{T'}(\mu, \nu)| &\geq \alpha.
%         \end{align*}
%     Therefore,
%         \begin{align*}
%             \delta(\mu, \nu) = \min\{|\wass_T(\mu, \nu) - \wass_{T'}(\mu, \nu)| : T, T' \in \setoftrees, \wass_T(\mu, \nu) \neq \wass_{T'}(\mu, \nu)\} \geq \alpha.
%         \end{align*}
% \end{proof}

\Cref{thm:steepest-descent-empirical} now follows immediately from~\cref{lem:rational-measures-gap} and~\cref{thm:steepest-descent-convergence}.

\section{Ollivier--Ricci curvature}\label{sec:ollivier-ricci-curvature}


\begin{thebibliography}{99}
\bibitem[bapat1997moore]{bapat1997moore} {\sc R.~B. Bapat}, {\em Moore-{P}enrose inverse of the incidence matrix of a tree}, Linear and Multilinear Algebra, 42 (1997), pp.~159--167.
\bibitem[essid2018quadratically]{essid2018quadratically} {\sc M.~Essid and J.~Solomon}, {\em Quadratically regularized optimal transport on graphs}, SIAM Journal on Scientific Computing, 40 (2018), pp.~A1961--A1986.
\bibitem[peyre2019computational]{peyre2019computational} {\sc G.~Peyr{\'e} and M.~Cuturi}, {\em Computational optimal transport}, Foundations and Trends in Machine Learning, 11 (2019), pp.~355--607.
\bibitem[robertson2024all]{robertson2024all} {\sc S.~Robertson, Z.~Wan, and A.~Cloninger}, {\em Resistance distance and linearized optimal transport on graphs}, arXiv preprint arXiv:2404.15261, (2024).
\bibitem[ziegler1995lectures]{ziegler1995lectures} {\sc G.~M. Ziegler}, {\em Lectures on polytopes}, vol.~152 of Graduate Texts in Mathematics, Springer-Verlag, New York, 1995.
\end{thebibliography}
\end{document}
